\chapter{Positive solutions and limits of normal graphs}\label{sec:jacobi-limits}
This section supplies the analytic steps behind separated-sheet convergence in \cite{ChMin}, Section 33, pp.~79--82. We retain Harnack and Schauder estimates as external inputs and prove the normalization, linearization and compactness deductions. The scalar mean-curvature convention is $H=\operatorname{div}_\Sigma\nu$, so its derivative at a minimal hypersurface is $-J$.

\begin{theorem}[Harnack with bounded lower-order coefficients]\label{thm:harnack-lower-order}
Let $n\geq1$, $R>0$, and let $A=(a^{ij})$ be symmetric and measurable on $B_{2R}\subset\mathbb R^n$, with $\lambda I\leq A\leq\Lambda I$ for $0<\lambda\leq\Lambda$. Let $b,c$ be bounded measurable coefficients. If $u\in H^1(B_{2R})$ is nonnegative and solves
\[
 \operatorname{div}(A Du)+b\cdot Du+cu=0
\]
weakly, then
\[
 \operatorname*{ess\,sup}_{B_R}u\leq C\operatorname*{ess\,inf}_{B_R}u,
\]
where $C$ depends on $n,\lambda,\Lambda,R\|b\|_\infty,R^2\|c\|_\infty$. No sign assumption on $c$ is required for this inequality for nonnegative solutions. The Harnack estimate is an external analytic input from \cite{ChMin}, Appendix B, p.~87, and \cite{SimonPDE}, Chapter 18. The nonnegativity and interior margin, implicit in the shorthand appendix statement, are explicit here.
\end{theorem}

\begin{lemma}[Compactness after normalizing a positive solution]\label{lem:normalized-positive-solutions}
Let $\Sigma$ be a connected smooth manifold of dimension $n\geq1$, exhausted by connected relatively compact open sets $U_j$, and fix $p\in U_1$. Let $L_j$ be smooth second-order operators on $U_j$ whose coefficients converge in $C^\infty$ on every compact subset to those of a locally uniformly elliptic operator $L$. Assume their principal coefficients are symmetric. If $v_j>0$, $L_jv_j=0$, and $v_j(p)=1$, then a subsequence converges in $C^\infty$ on compact subsets to $v>0$ with $Lv=0$ and $v(p)=1$.
\end{lemma}
\begin{proof}
On a coordinate chart write $L_j=a_j^{ab}\partial_{ab}+b_j^a\partial_a+c_j$. Rewriting its principal part as $\partial_a(a_j^{ab}\partial_b)$ changes the drift to $b_j^b-\partial_a a_j^{ab}$. Smooth coefficient convergence supplies uniform ellipticity and bounds for this drift and the potential on every fixed compact chart. Theorem \ref{thm:harnack-lower-order} therefore applies uniformly there.
Given a compact set $K\subset\Sigma$, use connectedness to join $p$ to finitely many coordinate balls covering $K$ by a finite chain of overlapping smaller balls inside a larger relatively compact open set. For all sufficiently large $j$ these balls lie in $U_j$. Iterating Harnack along the chain gives $C_K^{-1}\leq v_j\leq C_K$ on $K$. Interior Schauder estimates, including the higher estimates in Theorem \ref{thm:interior-schauder}, give uniform bounds for all derivatives on smaller charts. Arzela--Ascoli and a diagonal subsequence yield smooth local convergence. The positive lower bounds and normalization pass to the limit, and coefficient convergence gives $Lv=0$. Neither the Harnack iteration proof nor Schauder theory is being reproved here.
\end{proof}

\begin{lemma}[The equation for the gap between minimal normal graphs]\label{lem:normal-graph-gap-equation}
Let $\Sigma^n$ be a smooth two-sided minimal hypersurface in a smooth Riemannian manifold, with $n\geq1$. Suppose $u_{1,j},u_{2,j}$ define minimal normal graphs on an exhaustion and converge to zero smoothly on compact subsets. There are linear operators $L_j$, locally uniformly elliptic for large $j$, such that
\[
 L_j(u_{2,j}-u_{1,j})=0,\qquad L_j\longrightarrow J=\Delta_\Sigma+|B|^2+\operatorname{Ric}(\nu,\nu)
\]
in the smooth coefficient topology on compact subsets.
\end{lemma}
\begin{proof}
Pull scalar mean curvature of the graph back to $\Sigma$ and denote the resulting operator by $\mathcal H(u)$. In a normal coordinate tube $F(x,s)=\exp_x(s\nu)$, the induced metric and graph normal depend smoothly on $(x,u,Du)$. Tracing the second fundamental form shows that $\mathcal H$ is quasilinear of order two. More explicitly, its coefficient of $\partial_{ab}u$ is
\[
 -g_u^{ab}\langle\partial_sF,\nu_u\rangle.
\]
The inner product is positive for small $C^1$ graphs; hence the negative of the linearization is elliptic. The remaining coefficients of the linearization depend smoothly on $(x,u,Du,D^2u)$ and the ambient metric. For $w_j=u_{2,j}-u_{1,j}$, the fundamental theorem of calculus gives
\[
 0=\mathcal H(u_{2,j})-\mathcal H(u_{1,j})
   =\int_0^1 D\mathcal H(u_{1,j}+tw_j)[w_j],dt.
\]
Set $L_j=-\int_0^1D\mathcal H(u_{1,j}+tw_j)\,dt$. On each compact set all interpolating graphs tend smoothly to zero, so the coefficients converge smoothly to $-D\mathcal H(0)=J$ by Proposition \ref{prop:mean-curvature-linearization}. This also gives uniform local ellipticity. This is the operator Taylor expansion in \cite{ChMin}, Proposition 33.3, with its sign fixed by the stated scalar convention.
\end{proof}

\begin{corollary}[Comparable sheet gaps on compact subsets]\label{cor:sheet-gap-comparison}
Under the hypotheses of Theorem \ref{thm:separated-sheets-jacobi}, write $w_j=u_{2,j}-u_{1,j}>0$. For every compact $K$ in the interior there is $C_K\geq1$ such that, for all large $j$,
\[
 C_K^{-1}\leq w_j(x)/w_j(p)\leq C_K\qquad(x\in K).
\]
If a normalized subsequence tends to $v$ and $q$ is another interior point, normalization at $q$ instead gives the limit $v/v(q)$.
\end{corollary}
\begin{proof}
Apply the Harnack-chain step of Lemma \ref{lem:normalized-positive-solutions} to the operators from Lemma \ref{lem:normal-graph-gap-equation}. This gives the bounds without a priori control on the rate at which the unnormalized gap tends to zero. For the last assertion divide $w_j(x)/w_j(p)$ by $w_j(q)/w_j(p)$; the denominator converges to the positive number $v(q)$.
\end{proof}

\begin{proposition}[Positive supersolutions and the quadratic form]\label{prop:positive-supersolution-transform}
Let $(\Sigma,g)$ be smooth, let $q$ be smooth, and set $J=\Delta+q$ and $Q(\phi)=\int(|\nabla\phi|^2-q\phi^2)$. For every smooth $w>0$ and compactly supported smooth interior test $\phi$,
\[
 Q(\phi)=\int w^2|\nabla(\phi/w)|^2+\int\frac{-Jw}{w}\phi^2.
\]
In particular $Jw\leq0$ implies $Q\geq0$ on such tests. This extends the positive-solution identity used in \cite{ChMin}, Proposition 33.4.
\end{proposition}
\begin{proof}
Write $\phi=w\eta$. Expand $|\nabla(w\eta)|^2$ and group the terms other than $w^2|\nabla\eta|^2$ as $\nabla w\cdot\nabla(w\eta^2)-qw^2\eta^2$. Their integral is $-\int w\eta^2(\Delta w+qw)$ by integration by parts. Substitution of $\eta=\phi/w$ proves the identity. Compact support removes boundary terms, and $w$ has a positive minimum on the support, so every expression is well defined.
\end{proof}

\begin{proposition}[A positive Jacobi field on a closed manifold has simple kernel]\label{prop:positive-jacobi-simple-kernel}
Let $\Sigma$ be smooth, connected and closed, and let $q$ be smooth. If $w>0$ satisfies $(\Delta+q)w=0$, then the smooth kernel of $J=\Delta+q$ consists exactly of the constant multiples of $w$.
\end{proposition}
\begin{proof}
For any $Jf=0$, integration by parts gives $Q(f)=-\int fJf=0$. Proposition \ref{prop:positive-supersolution-transform}, applied on the compact manifold with no boundary, gives $0=\int w^2|\nabla(f/w)|^2$. Thus $f/w$ is constant by connectedness. Conversely every constant multiple of $w$ lies in the kernel. Closedness is essential to this integration argument; no corresponding assertion for all solutions on a noncompact manifold is made.
\end{proof}

\begin{proposition}[Accumulating minimal normal graphs force degeneracy]\label{prop:accumulating-graphs-jacobi}
Let $\Sigma^n$ be a smooth closed two-sided embedded minimal hypersurface, $n\geq1$. Suppose nonzero smooth functions $u_j\to0$ in $C^\infty(\Sigma)$ define minimal normal graphs. Then $J$ has a nonzero smooth kernel element. Consequently, if $\ker J=0$, the zero graph is isolated among minimal normal graphs in the smooth topology. The functions $u_j$ need not have a common sign. Compare the nullity mechanism in \cite{Sharp17}, Theorem 2.3 and its discussion.
\end{proposition}
\begin{proof}
Apply the averaged linearization of Lemma \ref{lem:normal-graph-gap-equation} with $u_{1,j}=0$ and $u_{2,j}=u_j$, and set $v_j=u_j/\|u_j\|_{C^0}$. Then $L_jv_j=0$, $L_j\to J$ smoothly and $\|v_j\|_{C^0}=1$. A finite covering of the closed hypersurface and interior Schauder estimates yield uniform bounds in every $C^k$ norm. A subsequence converges smoothly to $v$ with $Jv=0$ and $\|v\|_{C^0}=1$. If the zero graph were not isolated in the metrizable smooth topology, a sequence of distinct minimal graphs converging smoothly to zero would contradict the trivial-kernel assumption. This is local isolation of normal graphs, not a global finiteness theorem for minimal hypersurfaces.
\end{proof}
